\begin{table}[htbp]
\centering
\footnotesize
\caption{500,000 অর্ডারে রোপিত 1,000টি লেনদেন-ত্রুটির শনাক্তকরণ}
\label{tab:anomaly}
\begin{adjustbox}{max width=\linewidth}
\begin{tabular}{lrrrrrrr}
\toprule
\textbf{Method} & \textbf{Flagged} & \textbf{Precision} & \textbf{Recall} & \textbf{Price x10} & \textbf{Qty x20} & \textbf{Disc. 30-45\%} & \textbf{Price 10\%} \\
\midrule
IQR on log quantity (per product) & 1,500 & 0.112 & 0.168 & 0.00 & 0.66 & 0.00 & 0.01 \\
Business rules (price, discount) & 750 & 1.000 & 0.750 & 1.00 & 0.00 & 1.00 & 1.00 \\
Isolation Forest (4 features) & 2,500 & 0.364 & 0.909 & 1.00 & 0.64 & 1.00 & 1.00 \\
Rules OR Isolation Forest & 2,500 & 0.364 & 0.909 & 1.00 & 0.64 & 1.00 & 1.00 \\
\bottomrule
\end{tabular}
\end{adjustbox}
\par\smallskip\parbox{0.92\linewidth}{\footnotesize\emph{টীকা:} শেষ চারটি কলাম দেখায় প্রতিটি ধরনের ত্রুটির কত অংশ শনাক্ত হয়েছে। ডেটার একটি অনুলিপিতে প্রতিটি ধরনের 250টি ত্রুটি রোপণ করা হয়েছে।}
\end{table}